% SOURCE_AUTHORITY: sga5_fr_workpass.tex, lines 2311--2451 (LF-normalized unit).
% SOURCE_SHA256: 791F4EFFC5E02832D5D77ED03518C8156D6F07E4C8238B03545DB93D883FBB28
\subsection*{3.1}
Sean, para $i=1,2$, $X_i$ un $S$-esquema y $L_i\in\ob\Dctf(X_i)$,
\[
X=X_1\times_S X_2,
\qquad
p_i:X\longrightarrow X_i,
\qquad
q_i:X_i\longrightarrow S
\]
las proyecciones canónicas:
\[
\begin{tikzcd}[column sep=1.8em,row sep=1.5em]
& X \arrow[dl,"P_1"'] \arrow[dr,"P_2"] & \\
L_1\quad X_1 \arrow[dr,"q_1"'] & & X_2\quad L_2 \arrow[dl,"q_2"] \\
& S &
\end{tikzcd}
\]
Con las notaciones de 1.3, se tiene un isomorfismo canónico (2.2.4)
\[
DL_1\overset{\mathbf L}{\otimes}_S L_2
\xrightarrow{\sim}
\uRHom\bigl(L_1\overset{\mathbf L}{\otimes}_S A_{X_2},
q_1^!A_S\overset{\mathbf L}{\otimes}_S L_2\bigr).
\]
Al componerlo con el isomorfismo de Künneth
\[
q_1^!A_S\overset{\mathbf L}{\otimes}_S L_2
\xrightarrow{\sim}p_2^!L_2 \qquad (1.7.3),
\]
se obtiene un isomorfismo
\begin{equation}
DL_1\overset{\mathbf L}{\otimes}_S L_2
\xrightarrow{\sim}
\uRHom(p_1^*L_1,p_2^!L_2).
\tag{3.1.1}
\end{equation}
Este desempeñará un papel fundamental en todo lo que sigue. Para $X_i=S$, 3.1.1 se reduce a
\[
L_1^\vee\overset{\mathbf L}{\otimes}L_2\xrightarrow{\sim}\uRHom(L_1,L_2),
\]
caso particular de un isomorfismo válido para complejos perfectos sobre un topos anillado cualquiera (SGA 6 I 7.7). Pondremos
\begin{equation}
\uRHom(p_1^*L_1,p_2^!L_2)=\uRHom_S(L_1,L_2),
\qquad
\Hom(p_1^*L_1,p_2^!L_2)=\Hom_S(L_1,L_2).
\tag{3.1.2}
\end{equation}
Los elementos de $\Hom_S(L_1,L_2)$ se llamarán \emph{correspondencias cohomológicas de $L_1$ a $L_2$}. En lugar de ``correspondencia cohomológica de $A_{X_1}$ a $A_{X_2}$'', diremos simplemente \emph{correspondencia cohomológica de $X_1$ a $X_2$} (se sobreentiende: con coeficientes en $A$).

Si $X_1$ es \emph{liso}, de dimensión relativa pura $r_1$, lo mismo ocurre con $p_2$, y por el teorema de dualidad (SGA 4 XVIII 3.2.5), se tiene
\[
p_2^!L_2\xrightarrow{\sim}p_2^*L_2(r_1)[2r_1],
\]
por tanto
\begin{equation}
\begin{aligned}
&\uRHom_S(L_1,L_2)=\uRHom(p_1^*L_1,p_2^*L_2)(r_1)[2r_1],\\
&\Hom_S(L_1,L_2)=\H^{2r_1}\bigl(X,\uRHom(p_1^*L_1,p_2^*L_2)(r_1)\bigr).
\end{aligned}
\tag{3.1.3}
\end{equation}
En particular, el grupo de las correspondencias cohomológicas de $X_1$ a $X_2$ se identifica en este caso con $\H^{2r_1}(X,A(r_1))$.

\subsection*{3.2}
Las correspondencias cohomológicas que aparecen en la práctica están asociadas a $S$-morfismos $X_2\to X_1$, o, más generalmente, a ciclos de $X$, o a esquemas sobre $X$. Sea
\[
c:C\longrightarrow X
\]
un esquema sobre $X$ (a veces se dice que $c$ es una correspondencia entre $X_1$ y $X_2$), denotemos por $c_i=p_ic:C\to X_i$ las proyecciones. Llamaremos \emph{correspondencias cohomológicas de $L_1$ a $L_2$ con soporte en $c$} (o en $C$, si no hay riesgo de confusión) a los elementos del grupo
\[
\H^0\bigl(X,c_!\uRHom(c_1^*L_1,c_2^!L_2)\bigr)
\quad(=\Hom(c_1^*L_1,c_2^!L_2)\text{ si }c\text{ es propio}).
\]
Por la fórmula de inducción (SGA 4 XVIII 3.1.12.2), se tiene un isomorfismo canónico
\begin{equation}
c^!\uRHom(p_1^*L_1,p_2^!L_2)
\xrightarrow{\sim}
\uRHom(c_1^*L_1,c_2^!L_2),
\tag{3.2.1}
\end{equation}
de donde, al aplicar $c_!$ y componer con la flecha de adjunción $c_!c^!\to \Id$, se obtienen flechas canónicas
\begin{equation}
c_!\uRHom(c_1^*L_1,c_2^!L_2)
\longrightarrow
\uRHom(p_1^*L_1,p_2^!L_2),
\tag{3.2.2}
\end{equation}
\begin{equation}
\H^0\bigl(X,c_!\uRHom(c_1^*L_1,c_2^!L_2)\bigr)
\longrightarrow
\Hom_S(L_1,L_2).
\tag{3.2.3}
\end{equation}
Toda correspondencia de $L_1$ a $L_2$ con soporte en $C$ define, por tanto, canónicamente, mediante 3.2.3, una correspondencia (sin soporte) de $L_1$ a $L_2$.

\emph{Ejemplos.} a) Supongamos que $c$ es propio y que $c_2$ tiene dimensión Tor finita y dimensión relativa $\leq N$. El morfismo de traza (SGA $4^{1/2}$ Ciclo 2.3.3)
\[
\mathrm{Tr}_{c_2}:c_{2!}A_C\longrightarrow A_{X_2}(-N)[-2N]
\]
define por adjunción una correspondencia cohomológica de $A_{X_1}$ a $A_{X_2}(-N)[2N]$ con soporte en $C$, que llamaremos \emph{clase de $(c,N)$}, y denotaremos
\begin{equation}
\operatorname{cl}(c,N)\in \H^{-2N}\bigl(C,c_2^!A_{X_2}(-N)\bigr).
\tag{3.2.4}
\end{equation}
El caso más importante es aquel en que $N=0$, es decir, $c_2$ es finito (y de dimensión Tor finita), en cuyo caso $\operatorname{cl}(c)$ es una correspondencia cohomológica de $X_1$ a $X_2$ con soporte en $C$. Si $X_1$ es liso, de dimensión relativa pura $r_1$, se tiene, por 3.1.3 y 3.2.1,
\begin{equation}
\H^{-2N}\bigl(C,c_2^!A_{X_2}(-N)\bigr)
=
\H^{2(r_1-N)}\bigl(C,c^!A_X(r_1-N)\bigr).
\tag{3.2.5}
\end{equation}
Si además $c$ es una inmersión cerrada, el segundo miembro se reescribe
\[
\H^{2(r_1-N)}_C\bigl(X,A_X(r_1-N)\bigr),
\]
y la clase de $(c,N)$, considerada como elemento de este grupo, no es sino la clase de cohomología asociada a $c$ en el sentido de (SGA $4^{1/2}$ Ciclo 2.3.1), como se deduce inmediatamente de las definiciones.

b) Sea $F:X_2\to X_1$ un $S$-morfismo, denotemos por $c(F):C(F)\to X$ su grafo, es decir, la imagen de la sección de $p_2$ definida por $F$ (``el conjunto de los puntos $(Fx,x)$''): $c_2$ es un isomorfismo, y $c_1c_2^{-1}=F$. Se tienen, por tanto, isomorfismos canónicos
\begin{equation}
\uRHom\bigl(c_1(F)^*L_1,c_2(F)^!L_2\bigr)=\uRHom(F^*L_1,L_2),
\qquad
\Hom\bigl(c_1(F)^*L_1,c_2(F)^!L_2\bigr)=\Hom(F^*L_1,L_2).
\tag{3.2.6}
\end{equation}
En otros términos, las correspondencias cohomológicas de $L_1$ a $L_2$ con soporte en el grafo de $F$ son los ``levantamientos'' de $F$ a un homomorfismo $F^*L_1\to L_2$. He aquí dos ejemplos importantes de tales correspondencias:

(i) $X_1=X_2$, $L_1=L_2$, $F=1$, la identidad de $X_1$, cuyo grafo $\Delta$ es la diagonal. La correspondencia con soporte en $\Delta$ definida por la identidad de $L_1$ se llama \emph{correspondencia diagonal}. Es el elemento
\[
1\in \Hom(L_1,L_1)=\H^0\bigl(X,\Delta_*\Delta^!\uRHom_S(L_1,L_1)\bigr)
\]
(grupo que se identifica con $\H^{2r}_\Delta(X,\uRHom_S(L_1,L_1)(r))$ cuando $X_1$ es liso, de dimensión relativa pura $r$).

(ii) Supongamos que $S$ es el espectro de un cuerpo finito $\mathbb F_q$, que $X_1=X_2$, $L_1=L_2$, y tomemos como $F:X_1\to X_1$ el endomorfismo de Frobenius relativo a $\mathbb F_q$, identidad sobre el espacio subyacente y elevación a la potencia $q$-ésima sobre $\mathcal O_{X_1}$ (cf. XIV), y como levantamiento de $F$ a $L_1$ el isomorfismo canónico
\[
F^*L_1\xrightarrow{\sim}L_1
\]
(loc. cit. §2 prop. 1) (más exactamente, el iterado $a$-ésimo del isomorfismo de (loc. cit.), si $q=p^a$). La correspondencia cohomológica de $L_1$ a $L_1$ con soporte en el grafo de $F$ así definida se llama \emph{correspondencia de Frobenius} (relativa a $\mathbb F_q$).

c) Algunas correspondencias cohomológicas entre haces sobre curvas, relacionadas con los operadores de Hecke, han sido estudiadas por Langlands en ([13] §7).

d) Las correspondencias cohomológicas entre haces constantes (las más frecuentes en la naturaleza) proporcionan de manera natural, por ``integración'', correspondencias cohomológicas entre complejos, como veremos ahora.
